% Supplementary evaluation material for the ICSE 2027 submission.

\section{Additional Leakage-Safe Results}
\label{sec:appendix-leakage-safe}

Table~\ref{tab:appendix-rq1-audit} reports the complete Quo Vadis
leakage-safe comparison.  All methods use the same exact-sequence- and
source-group-disjoint partition.  EsCapturer-full preserves a useful ranking
signal (AUC $0.9168{\pm}0.0121$), but its lower Recall reduces F1 and several
frequency and graph baselines perform better.  We report this result to bound
the record-level benchmark rather than claim superiority under the audit.

\begin{table}[!t]
\caption{Quo Vadis leakage-safe results (mean $\pm$ standard deviation over
seeds 7/17/27).}
\label{tab:appendix-rq1-audit}
\centering
\scriptsize
\setlength{\tabcolsep}{4pt}
\begin{tabular}{lrr}
\hline
Method & F1 & AUC \\
\hline
EsCapturer-full & $0.7667{\pm}0.0191$ & $0.9168{\pm}0.0121$ \\
API-frequency RF & $0.8211{\pm}0.0050$ & $0.9487{\pm}0.0003$ \\
API $n$-gram SVM & $0.8231{\pm}0.0000$ & $0.9559{\pm}0.0000$ \\
GAT & $0.7833{\pm}0.0069$ & $0.9394{\pm}0.0065$ \\
R-GCN & $0.7890{\pm}0.0049$ & $0.9395{\pm}0.0056$ \\
\hline
\end{tabular}
\end{table}

\paragraph{Failure analysis.}
The audit test fold contains 741 benign and 254 malware records.  Across the
three EsCapturer-full runs, 66--83 malware records are missed, whereas only
30--35 benign records are false alarms.  Family-level Recall is consistently
high for Ransomware (0.922--0.931) but substantially lower for RAT
(0.250--0.333), Coinminer (0.250--0.542), and Trojan (0.410--0.615).  The
missed malware traces also have a median length of 133--154 calls, compared
with two calls for correctly detected malware.  These diagnostics suggest
that reduced unique-family coverage and long unseen traces, rather than
false-positive inflation, dominate the observed F1 loss.
